\providecommand{\Nb}{\mathcal{N}}
\providecommand{\GL}{\operatorname{GL}}
\providecommand{\Aut}{\operatorname{Aut}}
\providecommand{\tr}{\operatorname{tr}}
\providecommand{\HASH}{\operatorname{HASH}}
\providecommand{\deg}{\operatorname{deg}}

\chapter{Sesgos Geom\'etricos: Invarianza y Equivarianza}
\label{ch:priors}

Los cap\'itulos anteriores introdujeron el lenguaje matem\'atico del aprendizaje
profundo geom\'etrico (\emph{geometric deep learning}, GDL) ---grupos,
representaciones y los dominios geom\'etricos donde residen los datos--- y
argumentaron que el fracaso de las arquitecturas cl\'asicas sobre datos no
euclidianos es conceptual, y no meramente t\'ecnico. Este cap\'itulo convierte ese
diagn\'ostico en un principio constructivo. La tesis central es sencilla de
enunciar: un modelo debe respetar las simetr\'ias del problema que resuelve, y la
forma de garantizarlo es incorporar la simetr\'ia en la arquitectura en lugar de
confiar en que se aprenda de los datos. Las dos nociones que lo precisan son la
\emph{invarianza} y la \emph{equivarianza}. A partir de ellas derivamos un
esquema de dise\~no ---los \emph{sesgos geom\'etricos} del t\'itulo--- que organiza
los dominios y las simetr\'ias del GDL en torno a cinco temas recurrentes y
unifica las arquitecturas convolucionales, de atenci\'on y de grafos bajo una
sola construcci\'on geom\'etrica: el paso de mensajes.

\section{El paradigma del aprendizaje profundo geom\'etrico}
\label{sec:priors:paradigm}

El aprendizaje profundo geom\'etrico extiende el aprendizaje profundo m\'as all\'a
de los dominios euclidianos tradicionales ---rejillas regulares de p\'ixeles,
secuencias lineales de \emph{tokens}--- hacia estructuras geom\'etricas
arbitrarias. Las mol\'eculas son grafos tridimensionales con simetr\'ias
rotacionales; las redes sociales son grafos irregulares con propiedades
topol\'ogicas emergentes; las superficies f\'isicas son variedades riemannianas
con curvatura intr\'inseca. El GDL proporciona el marco matem\'atico necesario
para construir arquitecturas neuronales que respeten y exploten estas geometr\'ias
naturales en lugar de restringirse a representaciones sobre dominios euclidianos.

Una de las contribuciones fundamentales del GDL es la demostraci\'on de que las
arquitecturas aparentemente dispares del aprendizaje profundo ---las redes
neuronales convolucionales (\emph{convolutional neural networks}, CNN), los
\emph{transformers}, las redes de grafos--- son instancias particulares de un
principio unificador: el procesamiento equivariante de se\~nales sobre dominios
geom\'etricos con simetr\'ias espec\'ificas. Las CNN emergen como el caso de
rejillas regulares con invarianza traslacional; los \emph{transformers} como
grafos completos donde la atenci\'on implementa un paso de mensajes ponderado; y
las redes neuronales de grafos (\emph{graph neural networks}, GNN) como el caso
general de grafos arbitrarios con invarianza a permutaciones. Esta perspectiva
unificadora se apoya en herramientas matem\'aticas profundas de la teor\'ia de
grupos, la geometr\'ia diferencial y el an\'alisis arm\'onico, y proporciona no
solo fundamentos te\'oricos sobre por qu\'e funcionan las arquitecturas
existentes, sino tambi\'en principios de dise\~no para arquitecturas nuevas con
garant\'ias de estabilidad y generalizaci\'on.

\subsection{Las cinco Ges fundamentales}
\label{subsec:priors:five-gs}

El marco del GDL se estructura en torno a cinco conceptos fundamentales, las
``cinco Ges''~\citep{bronstein2017geometric, bronstein2021geometric},
que proporcionan un vocabulario unificado para entender y dise\~nar arquitecturas
neuronales que aprovechan la geometr\'ia donde residen los datos.

\begin{description}
\item[Geometr\'ia.] La estructura intr\'inseca del dominio de datos y sus
propiedades m\'etricas: variedades riemannianas (como esferas para datos
omnidireccionales), espacios hiperb\'olicos (para datos jer\'arquicos) y espacios
producto (para datos multiestructurados). La geometr\'ia determina qu\'e
operaciones son naturales y eficientes en cada dominio.
\item[Grupos.] Las simetr\'ias que preservan la estructura del problema bajo
transformaci\'on. Las m\'as comunes son las traslaciones $\R^d$ (CNN), las
rotaciones $\mathrm{SO}(n)$ (redes equivariantes rotacionales), las permutaciones
$S_n$ (redes de grafos) y los grupos de Lorentz (redes relativistas). Identificar
el grupo de simetr\'ia correcto es crucial para el dise\~no arquitect\'onico.
\item[Grafos.] Discretizaciones de dominios geom\'etricos continuos que capturan
relaciones locales de vecindad. Los grafos pueden ser regulares (mallas
uniformes), irregulares (redes sociales), dirigidos (redes causales) o din\'amicos
(sistemas evolutivos). Constituyen el sustrato computacional de muchas
arquitecturas de GDL.
\item[Geod\'esicas.] Caminos \'optimos que respetan la geometr\'ia intr\'inseca
del espacio. En variedades riemannianas las geod\'esicas minimizan la distancia;
en grafos corresponden a caminos m\'as cortos; en espacios de informaci\'on
definen direcciones de m\'aximo gradiente. Las geod\'esicas determinan c\'omo se
propaga la informaci\'on en las arquitecturas geom\'etricas.
\item[Gauges.] Invarianzas bajo transformaciones de \emph{gauge} (cambios locales
de marco de referencia) que preservan la f\'isica observable del sistema.
Incluyen invarianzas electromagn\'eticas, libertades de coordenadas en variedades
y elecciones de orientaci\'on en grafos no dirigidos. Los \emph{gauges} aseguran
que las representaciones aprendidas sean intr\'insecas al dominio geom\'etrico.
\end{description}

La \Cref{fig:priors:five-gs} dispone las cinco Ges en torno al paradigma central
del GDL, y la \Cref{fig:priors:hierarchy} complementa esta visi\'on exhibiendo su
estructura jer\'arquica, desde la capa m\'as fundamental (la geometr\'ia del
dominio) hasta el nivel m\'as espec\'ifico de abstracci\'on (los marcos locales).

\begin{figure}[h]
\centering
\begin{tikzpicture}[
    node distance=1.8cm,
    gnode/.style={
        circle, draw=gdlgray!80, fill=#1, minimum size=2.0cm,
        font=\bfseries\small
    },
    arrow/.style={-{Stealth[length=8pt]}, thick, draw=gdlgray!60},
    lbl/.style={font=\small\itshape, text=gdlgray!70}
]
\node[gnode=gdlgray!25, minimum size=2.4cm] (gdl) at (0,0) {GDL};
\node[gnode=gdlblue!20] (geometry) at (90:3.6cm) {Geometr\'ia};
\node[gnode=gdlgreen!20] (groups) at (162:3.6cm) {Grupos};
\node[gnode=gdlorange!20] (graphs) at (234:3.6cm) {Grafos};
\node[gnode=gdlred!20] (geodesics) at (306:3.6cm) {Geod\'esicas};
\node[gnode=gdlpurple!20] (gauges) at (18:3.6cm) {Gauges};
\draw[arrow] (gdl) -- (geometry);
\draw[arrow] (gdl) -- (groups);
\draw[arrow] (gdl) -- (graphs);
\draw[arrow] (gdl) -- (geodesics);
\draw[arrow] (gdl) -- (gauges);
\node[lbl, above=0.05cm of geometry] {dominios};
\node[lbl, left=0.05cm of groups] {simetr\'ias};
\node[lbl, below left=0.05cm of graphs] {discretizaci\'on};
\node[lbl, below right=0.05cm of geodesics] {distancias};
\node[lbl, right=0.05cm of gauges] {coordenadas};
\draw[thick, dashed, draw=gdlgray!40] (geometry) to[bend left=15] (groups);
\draw[thick, dashed, draw=gdlgray!40] (groups) to[bend left=15] (graphs);
\draw[thick, dashed, draw=gdlgray!40] (graphs) to[bend left=15] (geodesics);
\draw[thick, dashed, draw=gdlgray!40] (geodesics) to[bend left=15] (gauges);
\draw[thick, dashed, draw=gdlgray!40] (gauges) to[bend left=15] (geometry);
\end{tikzpicture}
\caption[Las cinco Ges del aprendizaje profundo geom\'etrico]{Las cinco Ges
fundamentales del aprendizaje profundo geom\'etrico. Cada componente representa
un aspecto esencial del marco unificador: \textbf{Geometr\'ia} (dominios),
\textbf{Grupos} (simetr\'ias), \textbf{Grafos} (discretizaci\'on),
\textbf{Geod\'esicas} (distancias) y \textbf{Gauges} (coordenadas locales).}
\label{fig:priors:five-gs}
\end{figure}

\begin{figure}[h]
\centering
\begin{tikzpicture}[scale=0.8,
    level/.style={rectangle, rounded corners=3pt, draw, minimum height=0.7cm,
        align=center, font=\small\bfseries},
    arrow/.style={->, thick, >=stealth, gray!70},
    desc/.style={font=\scriptsize, text=black!60, align=left}
]
\node[level, fill=gdlgreen!15, minimum width=10cm] (geom) at (0,0) {Geometr\'ia};
\node[level, fill=gdlgreen!25, minimum width=8cm] (grupos) at (0,1.2) {Grupos};
\node[level, fill=gdlgreen!35, minimum width=6cm] (grafos) at (0,2.4) {Grafos};
\node[level, fill=gdlgreen!45, minimum width=4cm] (geod) at (0,3.6) {Geod\'esicas};
\node[level, fill=gdlgreen!55, minimum width=2.5cm] (gauges) at (0,4.8) {Gauges};
\draw[arrow] (geom.north) ++(0,0.1) -- ++(0,0.3);
\draw[arrow] (grupos.north) ++(0,0.1) -- ++(0,0.3);
\draw[arrow] (grafos.north) ++(0,0.1) -- ++(0,0.3);
\draw[arrow] (geod.north) ++(0,0.1) -- ++(0,0.3);
\node[desc, anchor=west] at (geom.east) {\quad dominio $\Omega$ donde viven los datos};
\node[desc, anchor=west] at (grupos.east) {\quad simetr\'ias $\group$ que preservan estructura};
\node[desc, anchor=west] at (grafos.east) {\quad discretizaci\'on del dominio};
\node[desc, anchor=west] at (geod.east) {\quad propagaci\'on de informaci\'on};
\node[desc, anchor=west] at (gauges.east) {\quad invarianzas locales};
\node[font=\scriptsize, anchor=east] at (geom.west) {$\R^n$, $\mathcal{M}$, $\mathcal{G}$\quad};
\node[font=\scriptsize, anchor=east] at (grupos.west) {$\mathrm{SO}(n)$, $S_n$, $E(n)$\quad};
\node[font=\scriptsize, anchor=east] at (grafos.west) {$\mathcal{G}=(V,E)$\quad};
\node[font=\scriptsize, anchor=east] at (geod.west) {$d(x,y)$, paths\quad};
\node[font=\scriptsize, anchor=east] at (gauges.west) {frames\quad};
\end{tikzpicture}
\caption[Jerarqu\'ia de las cinco Ges]{Estructura jer\'arquica de las cinco Ges
del aprendizaje profundo geom\'etrico. Desde la base, la \textbf{Geometr\'ia}
define el dominio donde residen los datos; los \textbf{Grupos} capturan las
simetr\'ias del problema; los \textbf{Grafos} discretizan el dominio; las
\textbf{Geod\'esicas} determinan c\'omo se propaga la informaci\'on; y los
\textbf{Gauges} manejan las invarianzas locales en los sistemas de coordenadas.}
\label{fig:priors:hierarchy}
\end{figure}

Desde esta perspectiva unificadora, el GDL sostiene que el procesamiento efectivo
de se\~nales sobre dominios geom\'etricos requiere respetar tres principios
fundamentales.

\begin{description}
\item[Localidad geom\'etrica.] Las operaciones deben ser locales respecto a la
m\'etrica del dominio. Las arquitecturas geom\'etricas explotan la estructura
local mediante vecindarios definidos por la geometr\'ia subyacente, permitiendo
que la informaci\'on se propague de forma consistente con las distancias
intr\'insecas del espacio (formalizado mediante grafos en \Cref{ch:graphs}).
\item[Equivarianza.] Las transformaciones deben preservar las simetr\'ias
relevantes del problema. Una funci\'on equivariante conmuta con las acciones del
grupo de simetr\'ia, garantizando que transformaciones equivalentes de la entrada
produzcan transformaciones correspondientes en la salida
(\Cref{sec:priors:inv-equi}).
\item[Composicionalidad.] Las representaciones complejas se construyen
jer\'arquicamente componiendo operaciones m\'as simples. La composici\'on de capas
equivariantes preserva la equivarianza extremo a extremo
(\Cref{thm:priors:composition}), permitiendo arquitecturas profundas que
mantienen las propiedades geom\'etricas deseadas. Esta propiedad sustenta el
teorema de aproximaci\'on universal equivariante (\Cref{thm:priors:universal}),
que garantiza el poder expresivo de las redes equivariantes.
\end{description}

\section{Por qu\'e las redes tradicionales necesitan un sesgo geom\'etrico}
\label{sec:priors:motivation}

Una red neuronal artificial tradicional trata su entrada como un vector sin
estructura $x \in \R^n$. Una sola neurona calcula $y = \sigma(\inner{w}{x} + b)$,
combinando todas las $n$ coordenadas mediante una suma ponderada global.
Reinterpretado como un operador sobre una se\~nal en un dominio
$\Omega = \{1, \dots, n\}$, este dominio es un mero conjunto de \'indices: no
porta noci\'on alguna de \emph{localidad} (toda coordenada contribuye por igual,
sin concepto de cercan\'ia y lejan\'ia), ni de \emph{vecindario estructurado} (la
conexi\'on $w_1 x_1$ es estructuralmente id\'entica a $w_{100} x_{100}$), ni de
\emph{invarianza geom\'etrica} (permutar los \'indices generalmente cambia la
salida, aun cuando la permutaci\'on sea sem\'anticamente irrelevante). La red
ignora cualquier geometr\'ia impl\'icita en los datos.

Esta ceguera tiene un coste concreto. Sup\'ongase que la magnitud que deseamos
predecir no cambia bajo un grupo $\group$ de transformaciones ---rotaciones de
una imagen, reetiquetados de los nodos de un grafo---. Una red completamente
conectada no tiene raz\'on alguna, por construcci\'on, para producir la misma
respuesta ante $x$ y ante una entrada transformada $g \cdot x$. El remedio
habitual, el \emph{aumento de datos}, presenta a la red copias transformadas de
cada ejemplo; pero esto infla el conjunto de entrenamiento por un factor de
$\abs{\group}$, no ofrece garant\'ia de que la funci\'on aprendida sea realmente
invariante, y se degrada o fracasa por completo cuando $\group$ es grande o
continuo. Un sesgo geom\'etrico reemplaza esta esperanza por una garant\'ia: la
simetr\'ia se codifica en la estructura de la aplicaci\'on, de modo que se
cumple exactamente, para toda entrada, por construcci\'on.

\section{Funciones invariantes y equivariantes}
\label{sec:priors:inv-equi}

Las nociones de equivarianza e invarianza constituyen el n\'ucleo te\'orico del
GDL. Tras haber desarrollado el lenguaje de grupos y representaciones
(\Cref{ch:foundations}), estamos ahora en posici\'on de formalizar con precisi\'on
c\'omo una funci\'on puede respetar las simetr\'ias de un problema. Trabajamos con
un grupo $\group$ que act\'ua sobre los espacios relevantes. Recu\'erdese que una
\emph{acci\'on} de $\group$ sobre un conjunto $X$ es una aplicaci\'on
$\group \times X \to X$, escrita $(g, x) \mapsto g \cdot x$, compatible con la ley
de grupo: $e \cdot x = x$ y $g \cdot (h \cdot x) = (gh) \cdot x$. Cuando $X$ es un
espacio vectorial y cada $g$ act\'ua linealmente, la acci\'on es una
\emph{representaci\'on} $\rho_X \colon \group \to \GL(X)$. Escribimos $g \cdot x$
por $\rho_X(g)\, x$ en todo lo que sigue.

\begin{definition}[Funci\'on equivariante]
\label{def:priors:equivariance}
Sea $\group$ un grupo que act\'ua sobre $X$ v\'ia $\rho_X$ y sobre $Y$ v\'ia
$\rho_Y$. Una aplicaci\'on $\phi \colon X \to Y$ es \emph{$\group$-equivariante}
si
\[
  \phi(g \cdot x) = g \cdot \phi(x)
  \qquad \text{para todo } g \in \group,\ x \in X,
\]
donde $g \cdot x$ denota $\rho_X(g)\, x$ y $g \cdot \phi(x)$ denota
$\rho_Y(g)\, \phi(x)$; equivalentemente,
$\phi \circ \rho_X(g) = \rho_Y(g) \circ \phi$ para todo $g$.
\end{definition}

\begin{definition}[Funci\'on invariante]
\label{def:priors:invariance}
Una aplicaci\'on $\phi \colon X \to Y$ es \emph{$\group$-invariante} si
$\phi(g \cdot x) = \phi(x)$ para todo $g \in \group$ y $x \in X$. Esto equivale a
la $\group$-equivarianza cuando la acci\'on de $\group$ sobre $Y$ es trivial,
$\rho_Y(g) = \mathrm{id}$.
\end{definition}

La invarianza es, por tanto, el caso especial de equivarianza en el que $\group$
act\'ua trivialmente sobre la salida. La equivarianza es la noci\'on m\'as
fundamental: afirma que transformar la entrada produce una salida transformada de
forma correspondiente, de modo que la estructura geom\'etrica se \emph{preserva}
en lugar de descartarse. En muchas aplicaciones de GDL buscamos capas intermedias
equivariantes que preserven la estructura, seguidas de una capa final invariante.

La \Cref{fig:priors:equivariance} representa la \Cref{def:priors:equivariance}
como un cuadrado conmutativo: la equivarianza es exactamente la afirmaci\'on de
que los dos caminos de $X$ a $Y$ ---transformar y luego aplicar $\phi$, o aplicar
$\phi$ y luego transformar--- coinciden.

\begin{figure}[h]
\centering
\begin{tikzpicture}[
    node distance=2.6cm and 4cm,
    space/.style={circle, draw=gdlgray!80, fill=gdlblue!15, minimum size=1.5cm, font=\large},
    arrow/.style={-{Stealth[length=7pt]}, thick},
    lbl/.style={font=\small, midway}
]
\node[space] (X) {$X$};
\node[space, right=of X] (Y) {$Y$};
\node[space, below=of X] (gX) {$X$};
\node[space, below=of Y] (gY) {$Y$};
\draw[arrow, gdlblue!70!black] (X) -- node[lbl, above] {$\phi$} (Y);
\draw[arrow, gdlblue!70!black] (gX) -- node[lbl, below] {$\phi$} (gY);
\draw[arrow, gdlred!70!black] (X) -- node[left, align=center] {\small\itshape transformar\\[1pt] $\rho_X(g)$} (gX);
\draw[arrow, gdlred!70!black] (Y) -- node[lbl, right] {$\rho_Y(g)$} (gY);
\node[font=\Large] at ($(X)!0.5!(gY)$) {$\circlearrowleft$};
\node[font=\normalsize, below=0.4cm of gX, anchor=north west, xshift=0.6cm] {
    $\phi(\rho_X(g)\cdot x) = \rho_Y(g)\cdot \phi(x)$
};
\end{tikzpicture}
\caption[El cuadrado conmutativo de la equivarianza]{La equivarianza como
diagrama conmutativo. Una aplicaci\'on $\phi \colon X \to Y$ es
$\group$-equivariante cuando el cuadrado conmuta: transformar la entrada mediante
$\rho_X(g)$ y luego aplicar $\phi$ produce el mismo resultado que aplicar $\phi$
y luego transformar la salida mediante $\rho_Y(g)$.}
\label{fig:priors:equivariance}
\end{figure}

Las redes neuronales artificiales tradicionales carecen de equivarianza por
dise\~no. Para aproximar la invarianza requieren \emph{aumento de datos} con todas
las transformaciones del grupo $\group$, inflando el conjunto de entrenamiento
por un factor de $\abs{\group}$ sin garant\'ia de generalizaci\'on. El GDL
resuelve esto incorporando las simetr\'ias directamente en la arquitectura,
propagando la equivarianza a trav\'es de todas las capas.

\begin{example}[Tres simetr\'ias can\'onicas]
\label{ex:priors:symmetries}
La clasificaci\'on de im\'agenes es invariante a traslaciones: desplazar la imagen
de un gato no cambia la etiqueta, de modo que la funci\'on de puntuaci\'on debe
ser $\Z^2$-invariante. La segmentaci\'on sem\'antica, en cambio, es
\emph{equivariante} a traslaciones: desplazar la imagen desplaza la m\'ascara de
segmentaci\'on en la misma cantidad. Una funci\'on definida sobre los nodos de un
grafo debe ser equivariante al grupo de permutaciones $S_n$: reetiquetar los
nodos permuta las salidas de forma id\'entica, porque un grafo no posee un orden
can\'onico de sus v\'ertices. Cada tarea fija un grupo $\group$ y un
comportamiento requerido ---invariante o equivariante--- bajo \'el.
\end{example}

La raz\'on por la que la equivarianza es la primitiva adecuada para las
arquitecturas profundas es que se preserva bajo composici\'on.

\begin{theorem}[Composici\'on de aplicaciones equivariantes]
\label{thm:priors:composition}
Sea $\group$ un grupo que act\'ua sobre $X$, $Y$, $Z$ mediante las
representaciones $\rho_X, \rho_Y, \rho_Z$. Si $f \colon X \to Y$ y
$h \colon Y \to Z$ son $\group$-equivariantes, entonces tambi\'en lo es la
composici\'on $h \circ f \colon X \to Z$.
\end{theorem}

\begin{proof}
Fijemos $g \in \group$ y $x \in X$. Usando la $\group$-equivarianza de $f$ y
luego la de $h$,
\[
  (h \circ f)(g \cdot x) = h\bigl(f(g \cdot x)\bigr)
    = h\bigl(g \cdot f(x)\bigr)
    = g \cdot h\bigl(f(x)\bigr)
    = g \cdot (h \circ f)(x). \qedhere
\]
\end{proof}

La \Cref{thm:priors:composition} es la piedra angular de todo el marco. Garantiza
que una red construida apilando capas equivariantes es ella misma equivariante,
extremo a extremo ---la equivarianza es una propiedad que establecemos una vez,
por capa, y heredamos gratis en cualquier profundidad---. Componer ese tronco
equivariante con una sola lectura invariante al final produce un predictor
globalmente invariante. Esto convierte la simetr\'ia de una propiedad que
verificar \emph{a posteriori} en una receta constructiva para arquitecturas
completas.

\begin{figure}[h]
\centering
\begin{tikzpicture}[
    >=Stealth,
    layer/.style={rectangle, draw, thick, rounded corners, minimum width=1.9cm, minimum height=1.1cm, font=\small},
    data/.style={circle, draw, thick, minimum size=0.95cm, font=\small},
    arrow/.style={->, thick}
]
\node[data, fill=gdlblue!20] (x) at (0, 0) {$x$};
\node[layer, fill=gdlgreen!20] (f1) at (2.8, 0) {Capa $f_1$};
\node[layer, fill=gdlgreen!20] (f2) at (5.6, 0) {Capa $f_2$};
\node[layer, fill=gdlgreen!20] (f3) at (8.4, 0) {Capa $f_3$};
\node[data, fill=gdlred!20] (y) at (11.2, 0) {$y$};
\draw[arrow] (x) -- (f1);
\draw[arrow] (f1) -- (f2);
\draw[arrow] (f2) -- (f3);
\draw[arrow] (f3) -- (y);
\node[data, fill=gdlblue!40] (gx) at (0, -2.8) {$g \cdot x$};
\node[data, fill=gdlred!40] (gy) at (11.2, -2.8) {$g \cdot y$};
\draw[arrow, gdlpurple!70] (x) -- node[left, font=\small] {$\rho_X(g)$} (gx);
\draw[arrow, gdlpurple!70] (y) -- node[right, font=\small] {$\rho_Y(g)$} (gy);
\draw[arrow, dashed, gdlorange!70] (gx) -- (8.7, -2.8) -- (gy);
\node[font=\small, gdlpurple!70] at (1.4, -1.4) {transformar};
\node[font=\small, gdlorange!70] at (5.6, -3.3) {$(f_3 \circ f_2 \circ f_1)(g \cdot x)$};
\node[font=\Large] at (11.2, -1.4) {$=$};
\node[font=\tiny, above] at (2.8, 0.65) {$\rho_1$-equiv};
\node[font=\tiny, above] at (5.6, 0.65) {$\rho_2$-equiv};
\node[font=\tiny, above] at (8.4, 0.65) {$\rho_3$-equiv};
\end{tikzpicture}
\caption[Composici\'on de capas equivariantes]{Composici\'on de capas
equivariantes. Si cada capa $f_i$ es $\group$-equivariante, entonces la
composici\'on $f_3 \circ f_2 \circ f_1$ es $\group$-equivariante: transformar la
entrada y luego aplicar la red produce el mismo resultado que aplicar la red y
luego transformar la salida.}
\label{fig:priors:composition}
\end{figure}

\section{El esquema del aprendizaje profundo geom\'etrico}
\label{sec:priors:blueprint}

El principio de composici\'on sugiere una plantilla para dise\~nar arquitecturas
sobre cualquier dominio geom\'etrico~\citep{bronstein2021geometric}. Fijemos un
dominio $\Omega$ (una rejilla, un grafo, una variedad), un espacio de se\~nales
$\mathcal{X}(\Omega)$ sobre \'el y un grupo de simetr\'ia $\group$ que act\'ua
sobre ambos. El \emph{esquema del aprendizaje profundo geom\'etrico} ensambla un
modelo a partir de cuatro tipos de bloque constructivo.

\begin{description}
\item[Capas equivariantes.] Aplicaciones lineales o no lineales
$\mathcal{X}(\Omega) \to \mathcal{X}(\Omega')$ que conmutan con la acci\'on de
$\group$. Son el motor: extraen caracter\'isticas mientras se transforman de
forma predecible bajo la simetr\'ia. Una no linealidad puntual aplicada tras una
aplicaci\'on lineal equivariante preserva la equivarianza siempre que act\'ue de
forma id\'entica en cada posici\'on.
\item[Agrupamiento local (engrosamiento).] Aplicaciones equivariantes que reducen
la resoluci\'on del dominio, $\Omega \to \Omega'$, respetando la simetr\'ia. El
agrupamiento implementa la \emph{separaci\'on de escalas}: permite que las capas
posteriores vean vecindarios mayores a trav\'es de un dominio m\'as grueso,
construyendo una jerarqu\'ia de caracter\'isticas.
\item[Lectura global invariante.] Una aplicaci\'on $\group$-invariante final que
colapsa la se\~nal a la salida deseada ---una sola etiqueta, una propiedad a
nivel de grafo--- t\'ipicamente una agregaci\'on sim\'etrica (suma, media,
m\'aximo) seguida de una red est\'andar.
\item[Atenci\'on o funciones de mensaje que respetan la simetr\'ia.] Cuando la
conectividad misma debe adaptarse al contenido, los pesos que enlazan posiciones
se calculan mediante funciones que son a su vez equivariantes. Esto generaliza
los filtros locales fijos a filtros dependientes de los datos.
\end{description}

Dos principios estructurales sustentan esta plantilla y reaparecen a lo largo del
libro.

\begin{description}
\item[Simetr\'ia (equivarianza).] Toda operaci\'on interna conmuta con $\group$,
de modo que la funci\'on global es equivariante por la
\Cref{thm:priors:composition}, alcanz\'andose la invarianza solo en la lectura.
\item[Separaci\'on de escalas (localidad y composicionalidad).] Las operaciones
son \emph{locales} respecto a la geometr\'ia de $\Omega$ ---cada capa agrega solo
sobre un vecindario--- y las representaciones complejas se construyen de forma
\emph{composicional}, apilando capas locales intercaladas con engrosamiento. La
localidad mantiene tratables el n\'umero de par\'ametros y el c\'omputo; la
composicionalidad permite que una pila profunda de operaciones locales capture
estructura global de largo alcance.
\end{description}

\begin{remark}[Un esquema, cinco dominios]
\label{rem:priors:fivegs}
El esquema es deliberadamente independiente del dominio. Instanciarlo sobre una
rejilla con el grupo de traslaciones recupera las redes convolucionales
(\Cref{ch:grids}); extender el grupo a rotaciones y reflexiones da las redes
equivariantes a grupos (\Cref{ch:groups}); reemplazar la rejilla por un grafo
arbitrario con el grupo de permutaciones da las redes neuronales de grafos
(\Cref{ch:graphs}); dotar al dominio de distancias intr\'insecas da el
aprendizaje en variedades (\Cref{ch:geodesics}); y reconocer la ausencia de un
marco local can\'onico en un dominio curvo da las redes equivariantes \emph{gauge}
(\Cref{ch:gauges}). Las cinco Ges son cinco realizaciones concretas de una sola
idea geom\'etrica.
\end{remark}

\section{Universalidad y expresividad}
\label{sec:priors:universality}

Una preocupaci\'on natural es que restringir una red a ser equivariante pueda
mermar su poder expresivo: la restricci\'on recorta un subespacio estricto de
todas las funciones. El siguiente resultado, que extiende el teorema cl\'asico de
aproximaci\'on universal al caso equivariante, disipa esa preocupaci\'on. Se debe
a \citet{yarotsky2022} y \citet{kondor2018generalizationconvolutionalneural}.

\begin{theorem}[Aproximaci\'on universal equivariante]
\label{thm:priors:universal}
Sea $\group$ un grupo compacto que act\'ua continuamente sobre espacios m\'etricos
compactos $X$ e $Y$, y sea $\mathcal{F}^{\group}(X, Y)$ el espacio de todas las
aplicaciones continuas $\group$-equivariantes $f \colon X \to Y$. Entonces, para
toda $f \in \mathcal{F}^{\group}(X, Y)$ y todo $\varepsilon > 0$ existe una red
neuronal $\group$-equivariante $\mathcal{N}$ tal que
\[
  \sup_{x \in X} d_Y\bigl(f(x), \mathcal{N}(x)\bigr) < \varepsilon,
\]
donde $d_Y$ es la m\'etrica en $Y$.
\end{theorem}

\begin{proof}
La demostraci\'on procede en cuatro pasos, siguiendo las ideas de
\citet{yarotsky2022}.

\emph{Paso 1: aproximaci\'on no equivariante.} Por el teorema cl\'asico de
aproximaci\'on universal, para cualquier $f \colon X \to Y$ continua y cualquier
$\varepsilon > 0$ existe una red neuronal sin restricciones
$\tilde f \colon X \to Y$ con
$\sup_{x \in X} d_Y(f(x), \tilde f(x)) < \varepsilon$.

\emph{Paso 2: promediado equivariante.} Definimos el \emph{operador de Reynolds}
promediando sobre el grupo respecto a la medida de Haar normalizada $\mu$ sobre
$\group$:
\[
  \bar f(x) = \int_{\group} \rho_Y(g^{-1})\, \tilde f(g \cdot x)\, d\mu(g),
\]
donde $\rho_Y \colon \group \to \GL(Y)$ es la representaci\'on de $\group$ en $Y$.

\emph{Paso 3: verificaci\'on de la equivarianza.} Comprobamos que $\bar f$ es
$\group$-equivariante. Para cualquier $h \in \group$,
\[
  \bar f(h \cdot x)
    = \int_{\group} \rho_Y(g^{-1})\, \tilde f(g \cdot (h \cdot x))\, d\mu(g)
    = \int_{\group} \rho_Y(g^{-1})\, \tilde f((gh) \cdot x)\, d\mu(g).
\]
Cambiando de variable a $g' = gh$ y usando la invarianza de la medida de Haar
($d\mu(g) = d\mu(g')$),
\[
  \bar f(h \cdot x)
    = \int_{\group} \rho_Y\bigl((g'h^{-1})^{-1}\bigr)\, \tilde f(g' \cdot x)\, d\mu(g')
    = \rho_Y(h) \int_{\group} \rho_Y(g'^{-1})\, \tilde f(g' \cdot x)\, d\mu(g')
    = \rho_Y(h)\, \bar f(x).
\]
Por tanto $\bar f$ es $\group$-equivariante.

\emph{Paso 4: control del error de aproximaci\'on.} Sea $f$ la funci\'on objetivo
$\group$-equivariante. Por la linealidad de la integral y la equivarianza de $f$,
\[
  d_Y\bigl(\bar f(x), f(x)\bigr)
    \le \int_{\group} d_Y\bigl(\rho_Y(g^{-1}) \tilde f(g \cdot x),\,
        \rho_Y(g^{-1}) f(g \cdot x)\bigr)\, d\mu(g).
\]
Para un grupo compacto las representaciones pueden tomarse unitarias, de modo que
$\rho_Y(g^{-1})$ preserva distancias y
\[
  d_Y\bigl(\bar f(x), f(x)\bigr)
    \le \int_{\group} d_Y\bigl(\tilde f(g \cdot x), f(g \cdot x)\bigr)\, d\mu(g)
    \le \sup_{x' \in X} d_Y\bigl(\tilde f(x'), f(x')\bigr) < \varepsilon,
\]
donde el \'ultimo paso usa que la acci\'on de $\group$ sobre el compacto $X$
preserva la compacidad. Aproximar $\bar f$ por una red $\group$-equivariante
completa la demostraci\'on. \qedhere
\end{proof}

El mensaje es decisivo: \emph{la equivarianza no cuesta nada en expresividad}. Una
red $\group$-equivariante puede aproximar cualquier objetivo continuo
$\group$-equivariante con precisi\'on arbitraria. La restricci\'on de simetr\'ia no
encoge el conjunto de funciones alcanzables dentro del subespacio relevante;
act\'ua \'unicamente como un sesgo inductivo, descartando las funciones no
equivariantes irrelevantes y mejorando as\'i la eficiencia muestral y la
generalizaci\'on sin sacrificar capacidad.

\begin{example}[Las CNN como caso especial]
\label{ex:priors:cnn-universal}
Para el grupo de traslaciones $\group = \Z^2$ actuando sobre im\'agenes discretas,
el teorema garantiza que las CNN ---que son equivariantes a traslaciones
(\Cref{ch:grids})--- pueden aproximar cualquier funci\'on continua invariante o
equivariante a traslaciones. Esto justifica te\'oricamente el \'exito emp\'irico
de las redes convolucionales en visi\'on por computador~\citep{cohen2016group}.
\end{example}

\section{El paso de mensajes: el marco unificador}
\label{sec:priors:message-passing}

Las cinco Ges forman un marco conceptual profundamente interconectado: la
\emph{geometr\'ia} determina la estructura del dominio y qu\'e \emph{grupos} de
simetr\'ias son naturales; los \emph{grafos} discretizan las geometr\'ias
continuas preservando propiedades topol\'ogicas y espectrales; las
\emph{geod\'esicas} definen caminos \'optimos que respetan la geometr\'ia; y los
\emph{gauges} aseguran que las representaciones sean intr\'insecas. Desde esta
perspectiva, las diferencias entre arquitecturas neuronales reflejan elecciones
sobre tres aspectos geom\'etricos: la \emph{estructura del dominio} (espacios
euclidianos, grafos, variedades), el \emph{grupo de simetr\'ias} (trivial,
traslaciones, permutaciones) y la \emph{conectividad} (global, vecindarios fijos,
adaptativa). Las redes neuronales artificiales tradicionales, operando en $\R^d$
con simetr\'ias triviales, constituyen el caso base: no explotan estructura
geom\'etrica alguna, de modo que toda invarianza debe aprenderse de los datos. El
GDL resuelve esto incorporando las simetr\'ias directamente en la arquitectura
mediante el principio de equivarianza.

El esquema se vuelve algor\'itmicamente concreto a trav\'es de una sola primitiva
computacional: el \emph{paso de mensajes}. Viendo los datos como un grafo
$\mathcal{G} = (V, E)$ de nodos y aristas, una capa propaga informaci\'on haciendo
que cada nodo agregue caracter\'isticas transformadas de sus vecinos. Esta sola
construcci\'on subsume las arquitecturas principales del aprendizaje profundo; sus
diferencias aparentes se reducen a elecciones de topolog\'ia de grafo,
conectividad y grupo de simetr\'ia.

\begin{definition}[Red neuronal de paso de mensajes]
\label{def:priors:mpnn}
Una \emph{red neuronal de paso de mensajes} (\emph{message-passing neural
network}, MPNN) opera sobre un grafo $\mathcal{G} = (V, E)$ con caracter\'isticas
de nodo $\{h_v\}_{v \in V}$ y de arista $\{e_{uv}\}_{(u,v) \in E}$. La
actualizaci\'on en la capa $\ell$ consta de dos fases, una de \emph{paso de
mensajes} y otra de \emph{actualizaci\'on de nodo}:
\[
  m_v^{(\ell+1)} = \bigoplus_{u \in \Nb(v)}
      M_\ell\bigl(h_v^{(\ell)}, h_u^{(\ell)}, e_{vu}\bigr),
  \qquad
  h_v^{(\ell+1)} = U_\ell\bigl(h_v^{(\ell)}, m_v^{(\ell+1)}\bigr),
\]
donde $\Nb(v)$ es el vecindario de $v$, $M_\ell$ una funci\'on de mensaje
aprendible, $\bigoplus$ un agregador permutaci\'on-invariante (suma, media,
m\'aximo) y $U_\ell$ una funci\'on de actualizaci\'on~\citep{gilmer2017neural}.
Este esquema opera exclusivamente sobre nodos y aristas; su generalizaci\'on a
$k$-s\'implices de orden superior, el \emph{paso de mensajes simplicial}, se
desarrolla en la \Cref{sec:priors:tdl}.
\end{definition}

El agregador $\bigoplus$ es permutaci\'on-invariante por dise\~no, lo que hace cada
capa equivariante a reetiquetados de los vecinos ---la capa de paso de mensajes es
la capa equivariante del esquema para datos relacionales---. Las siguientes
correspondencias (\Cref{fig:priors:mpnn}) muestran que las arquitecturas
familiares son casos especiales.

\begin{figure}[h]
\centering
\begin{tikzpicture}[>=Stealth, font=\small,
    nd/.style={circle, draw, thick, fill=white, minimum size=4mm, inner sep=0pt},
    edge/.style={thick, draw=gdlgray!60},
    box/.style={rectangle, rounded corners=4pt, draw=gdlgray!70, thick, minimum width=4.6cm, minimum height=3.7cm},
    title/.style={font=\small\bfseries, fill=white, inner sep=2pt},
    info/.style={font=\scriptsize, text=black!70, align=center}
]
% ANN
\begin{scope}[shift={(-3.3, 2.7)}]
    \node[box, fill=gdlblue!15] (annbox) {};
    \node[title] at (0, 2.05) {ANN (completamente conectada)};
    \foreach \i/\a in {1/90,2/162,3/234,4/306,5/18} \node[nd, fill=gdlblue!20] (a\i) at (\a:0.9) {};
    \foreach \i in {1,...,5} \foreach \j in {1,...,5} {\ifnum\i<\j \draw[edge, gdlblue!30] (a\i)--(a\j);\fi}
    \node[info] at (0, -1.45) {grafo completo $K_d$\\pesos espec\'ificos $W_{ij}$};
\end{scope}
% CNN
\begin{scope}[shift={(3.3, 2.7)}]
    \node[box, fill=gdlgreen!15] (cnnbox) {};
    \node[title] at (0, 2.05) {CNN (convoluci\'on)};
    \foreach \i in {0,1,2} \foreach \j in {0,1,2} \node[nd, fill=gdlgreen!20] (c\i\j) at ({\i*0.7-0.7},{\j*0.7-0.7}) {};
    \foreach \i in {0,1,2} \foreach \j [evaluate=\j as \jn using int(\j+1)] in {0,1} {
        \draw[edge, gdlgreen!50] (c\i\j)--(c\i\jn); \draw[edge, gdlgreen!50] (c\j\i)--(c\jn\i);}
    \draw[thick, gdlorange, rounded corners=2pt] (-0.3,-0.3) rectangle (0.3,0.3);
    \node[info] at (0, -1.45) {rejilla $\Z^d$, pesos compartidos\\$\group=(\Z^d,+)$};
\end{scope}
% Transformer
\begin{scope}[shift={(-3.3, -2.7)}]
    \node[box, fill=gdlpurple!15] (transbox) {};
    \node[title] at (0, 2.05) {Transformer (atenci\'on)};
    \foreach \i/\a in {1/90,2/162,3/234,4/306,5/18} \node[nd, fill=gdlpurple!20] (t\i) at (\a:0.9) {};
    \draw[edge, gdlpurple!60, line width=2pt] (t1)--(t3);
    \draw[edge, gdlpurple!40, line width=1.4pt] (t1)--(t2);
    \draw[edge, gdlpurple!20, line width=0.5pt] (t1)--(t4);
    \draw[edge, gdlpurple!50, line width=1.4pt] (t2)--(t4);
    \draw[edge, gdlpurple!60, line width=2pt] (t3)--(t5);
    \draw[edge, gdlpurple!30, line width=0.8pt] (t4)--(t5);
    \node[info] at (0, -1.45) {grafo completo $K_n$\\pesos adaptativos $\alpha_{ij}$};
\end{scope}
% GNN
\begin{scope}[shift={(3.3, -2.7)}]
    \node[box, fill=gdlred!15] (gnnbox) {};
    \node[title] at (0, 2.05) {GNN (red de grafos)};
    \node[nd, fill=gdlred!20] (g1) at (-0.7,0.6) {};
    \node[nd, fill=gdlred!20] (g2) at (0.3,0.8) {};
    \node[nd, fill=gdlred!20] (g3) at (0.8,0.1) {};
    \node[nd, fill=gdlred!20] (g4) at (-0.4,-0.3) {};
    \node[nd, fill=gdlred!20] (g5) at (0.4,-0.5) {};
    \node[nd, fill=gdlred!20] (g6) at (-0.15,0.25) {};
    \foreach \a/\b in {g1/g2,g1/g4,g1/g6,g2/g3,g2/g6,g3/g5,g4/g5,g4/g6,g5/g6}
        \draw[edge, gdlred!50] (\a)--(\b);
    \node[info] at (0, -1.45) {grafo arbitrario $\mathcal{G}$\\$\group=\Aut(\mathcal{G})$};
\end{scope}
% Centro
\node[font=\bfseries, align=center, fill=gdlyellow!20, rounded corners=3pt,
      draw=gdlyellow!70, thick, inner sep=5pt] (mpnn) at (0,0)
      {MPNN\\[-1pt]{\scriptsize $m_v=\bigoplus_{u\in\Nb(v)} M(h_v,h_u,e_{vu})$}};
\draw[->, thick, gdlblue!60] (annbox.south east) -- (mpnn.north west);
\draw[->, thick, gdlgreen!50!black] (cnnbox.south west) -- (mpnn.north east);
\draw[->, thick, gdlpurple!60] (transbox.north east) -- (mpnn.south west);
\draw[->, thick, gdlred!60] (gnnbox.north west) -- (mpnn.south east);
\end{tikzpicture}
\caption[Arquitecturas como paso de mensajes]{Las arquitecturas principales del
aprendizaje profundo como instancias del paso de mensajes. Difieren solo en la
topolog\'ia del grafo, el esquema de compartici\'on de pesos y el grupo de
simetr\'ia $\group$: redes completamente conectadas (grafo completo, pesos
espec\'ificos por par), redes convolucionales (rejilla, pesos compartidos por
traslaci\'on), \emph{transformers} (grafo completo, pesos adaptativos al
contenido) y redes de grafos (grafo arbitrario). La actualizaci\'on de paso de
mensajes del centro es com\'un a todas.}
\label{fig:priors:mpnn}
\end{figure}

\subsection{Las capas completamente conectadas como MPNN triviales}
\label{subsec:priors:ann-mpnn}

\begin{proposition}[ANN como MPNN sobre un grafo completo uniforme]
\label{prop:priors:ann-mpnn}
Una capa completamente conectada $h^{(\ell+1)} = \sigma(W h^{(\ell)} + b)$ es
equivalente a una MPNN donde el grafo es el grafo completo $K_d$ cuyas
caracter\'isticas son nodos, el mensaje es el espec\'ifico por par
$M(h_v, h_u, e_{vu}) = W_{vu} h_u$, la agregaci\'on es la suma sobre todas las
caracter\'isticas de entrada y la actualizaci\'on es
$U(h_v, m_v) = \sigma(m_v + b_v)$.
\end{proposition}

\begin{proof}
Sea $h^{(\ell)} \in \R^d$ el vector de caracter\'isticas en la capa $\ell$. La
capa completamente conectada calcula
\[
  h_v^{(\ell+1)} = \sigma\!\left(\sum_{u=1}^{d} W_{vu}\, h_u^{(\ell)} + b_v\right).
\]
Interpretando cada componente como un nodo del grafo completo $K_d$, esta
expresi\'on coincide exactamente con el paso de mensajes de
\Cref{def:priors:mpnn}, con mensaje
$M(h_v, h_u, e_{vu}) = W_{vu} h_u^{(\ell)}$, agregaci\'on
$m_v = \sum_{u \in \Nb(v)} M(\cdot) = \sum_{u=1}^d W_{vu} h_u^{(\ell)}$ y
actualizaci\'on $U(h_v, m_v) = \sigma(m_v + b_v)$. \qedhere
\end{proof}

Esto revela por qu\'e las redes completamente conectadas carecen de \emph{sesgo
inductivo geom\'etrico}: tratan toda conexi\'on entre caracter\'isticas como
igualmente importante, ignorando cualquier estructura espacial, topol\'ogica o de
similaridad que los datos puedan poseer.

\subsection{Las convoluciones como paso de mensajes regular}
\label{subsec:priors:cnn-mpnn}

Las CNN son la primera instancia no trivial del paso de mensajes, en la que la
estructura de rejilla impone un sesgo inductivo geom\'etrico fundamental: la
equivarianza traslacional.

\begin{proposition}[CNN como MPNN sobre una rejilla regular]
\label{prop:priors:cnn-mpnn}
Una capa convolucional sobre una rejilla regular $\Z^d$ con soporte de n\'ucleo
$\mathcal{K}$ es equivalente a una MPNN con vecindario
$\Nb(v) = \{u \in \Z^d : v - u \in \mathcal{K}\}$, mensaje
$M(h_v, h_u, e_{vu}) = W_{v-u}\, h_u$ (pesos compartidos por posici\'on
relativa), agregaci\'on por suma sobre el vecindario local y actualizaci\'on
$U(h_v, m_v) = \sigma(m_v + b)$.
\end{proposition}

\begin{proof}
Sea $h^{(\ell)} \colon \Z^d \to \R^{c_{\mathrm{in}}}$ la se\~nal en la capa
$\ell$ y $\{W_\delta\}_{\delta \in \mathcal{K}}$ el n\'ucleo convolucional con
$W_\delta \in \R^{c_{\mathrm{out}} \times c_{\mathrm{in}}}$. La convoluci\'on
discreta calcula
\[
  h_v^{(\ell+1)}
    = \sigma\!\left(\sum_{\delta \in \mathcal{K}} W_\delta\, h_{v-\delta}^{(\ell)} + b\right)
    = \sigma\!\left(\sum_{u \in \Nb(v)} W_{v-u}\, h_u^{(\ell)} + b\right),
\]
donde el cambio de variable $u = v - \delta$ transforma la suma sobre
desplazamientos del n\'ucleo en una suma sobre los vecinos de $v$ en el grafo de
la rejilla. Esto coincide con el paso de mensajes de \Cref{def:priors:mpnn} con
mensaje $M(h_v, h_u, e_{vu}) = W_{v-u} h_u^{(\ell)}$, agregaci\'on
$m_v = \sum_{u \in \Nb(v)} M(\cdot)$ y actualizaci\'on
$U(h_v, m_v) = \sigma(m_v + b)$. \qedhere
\end{proof}

La diferencia clave respecto a las redes completamente conectadas es que los pesos
$W_{v-u}$ dependen solo de la posici\'on relativa, no del par absoluto $(v, u)$.
Esta \emph{compartici\'on de pesos} es precisamente la implementaci\'on de la
equivarianza traslacional como sesgo inductivo geom\'etrico, reduciendo
dr\'asticamente el n\'umero de par\'ametros de $O(n^2 c^2)$ a
$O(\abs{\mathcal{K}}\, c^2)$.

\subsection{La atenci\'on como paso de mensajes adaptativo}
\label{subsec:priors:attention-mpnn}

Los \emph{transformers} representan el caso complementario: los pesos de
comunicaci\'on se calculan din\'amicamente a partir del contenido de los nodos, no
de su posici\'on.

\begin{proposition}[Transformer como MPNN sobre un grafo completo adaptativo]
\label{prop:priors:attention-mpnn}
Una capa de autoatenci\'on sobre una secuencia de $n$ \emph{tokens} es
equivalente a una MPNN sobre el grafo completo $K_n$ cuyos nodos son los
\emph{tokens}, con mensaje $M(h_i, h_j) = \alpha_{ij}\, V h_j$ y coeficientes de
atenci\'on
\[
  \alpha_{ij}
    = \frac{\exp\!\bigl((Q h_i)^{\!\top}(K h_j)/\sqrt{d_k}\bigr)}
           {\sum_{m=1}^{n} \exp\!\bigl((Q h_i)^{\!\top}(K h_m)/\sqrt{d_k}\bigr)},
\]
agregaci\'on por suma ponderada $m_i = \sum_{j=1}^{n} M(h_i, h_j)$ y
actualizaci\'on identidad $U(h_i, m_i) = m_i$.
\end{proposition}

\begin{proof}
Sea $\{h_1, \dots, h_n\} \subset \R^d$ la secuencia de entrada y
$Q, K, V \in \R^{d \times d_k}$ las matrices de proyecci\'on. La capa de
autoatenci\'on calcula
\[
  h_i^{(\ell+1)} = \sum_{j=1}^{n} \alpha_{ij}\, V h_j^{(\ell)}
    = \sum_{j=1}^{n} \underbrace{\alpha_{ij}\, V h_j^{(\ell)}}_{M(h_i, h_j)}.
\]
Identificando el grafo completo $K_n$ con $\Nb(i) = \{1, \dots, n\}$, esto es
exactamente el paso de mensajes de \Cref{def:priors:mpnn} con mensaje
$M(h_i, h_j) = \alpha_{ij} V h_j$ y actualizaci\'on identidad
$U(h_i, m_i) = m_i$. La normalizaci\'on \emph{softmax} garantiza
$\alpha_{ij} \in (0, 1)$ con $\sum_{j=1}^{n} \alpha_{ij} = 1$, de modo que la
agregaci\'on es una combinaci\'on convexa, a diferencia de la suma uniforme de las
capas completamente conectadas (\Cref{prop:priors:ann-mpnn}) y las convoluciones
(\Cref{prop:priors:cnn-mpnn})~\citep{vaswani2017attentionneed}. \qedhere
\end{proof}

La diferencia fundamental respecto a las arquitecturas anteriores es que los pesos
$\alpha_{ij}$ son funciones del contenido $(h_i, h_j)$, ni fijos (completamente
conectadas) ni dependientes solo de la posici\'on relativa (convoluciones). Esta
adaptabilidad permite conectividad global a costa de una complejidad cuadr\'atica
$O(n^2)$ en la longitud de secuencia.

Leer en horizontal estos casos revela un \'unico eje de sofisticaci\'on
geom\'etrica creciente: las redes completamente conectadas no imponen sesgo
alguno (pesos uniformes espec\'ificos por par); las convoluciones imponen un sesgo
\emph{espacial} (pesos compartidos por posici\'on relativa, que codifican la
equivarianza traslacional); los \emph{transformers} reemplazan los pesos fijos por
pesos \emph{adaptativos, dependientes del contenido} (conectividad global a coste
cuadr\'atico); y las redes de grafos generalizan todo el esquema a topolog\'ias
relacionales arbitrarias, equivariantes al grupo de automorfismos
$\Aut(\mathcal{G})$. El sesgo geom\'etrico es precisamente lo que cada
arquitectura supone sobre qu\'e conexiones importan y c\'omo se comparten sus
pesos.

\subsection{El aprendizaje profundo topol\'ogico como paso de mensajes de orden superior}
\label{sec:priors:tdl}

Las tres arquitecturas anteriores (ANN, CNN, \emph{transformers}) operan
exclusivamente sobre los \emph{nodos} ($0$-s\'implices) de un grafo, propagando
informaci\'on a lo largo de aristas. El aprendizaje profundo topol\'ogico
(\emph{topological deep learning}, TDL) rompe esta restricci\'on: asigna se\~nales
a $k$-s\'implices de orden arbitrario y propaga informaci\'on simult\'aneamente a
trav\'es de las adyacencias inferior y superior definidas por los operadores de
frontera. Esta extensi\'on incorpora un \emph{sesgo inductivo topol\'ogico} que
captura relaciones de orden superior inaccesibles al paso de mensajes est\'andar.

\begin{definition}[Se\~nal simplicial]
\label{def:priors:simplicial-signal}
Una \emph{se\~nal simplicial de orden $k$} sobre un complejo simplicial
$\mathcal{K}$ es una funci\'on $\mathbf{x} \colon \mathcal{K}_k \to \R^d$ que
asigna un vector de caracter\'isticas $\mathbf{x}_\sigma \in \R^d$ a cada
$k$-s\'implice $\sigma \in \mathcal{K}_k$.
\end{definition}

\begin{definition}[Paso de mensajes simplicial]
\label{def:priors:simplicial-mp}
El \emph{paso de mensajes simplicial} generaliza el paso de mensajes en grafos
(\Cref{def:priors:mpnn}) mediante tres tipos de comunicaci\'on para cada
$k$-s\'implice $\sigma$:
\begin{enumerate}
\item \textbf{Mensaje inferior} (v\'ia $\partial_k^{\!\top}$): informaci\'on de
los $(k{-}1)$-s\'implices en la frontera de $\sigma$,
\[
  \mathbf{m}_\sigma^{\mathrm{down}}
    = \bigoplus_{\tau \in \mathcal{B}^{\mathrm{down}}(\sigma)}
      M^{\mathrm{down}}\bigl(\mathbf{h}_\sigma, \mathbf{h}_\tau\bigr),
\]
donde $\mathcal{B}^{\mathrm{down}}(\sigma) =
\{\tau : [\partial_k^{\!\top} \partial_k]_{\sigma,\tau} \neq 0\}$ son los
$k$-s\'implices adyacentes por caras compartidas.
\item \textbf{Mensaje superior} (v\'ia $\partial_{k+1}$): informaci\'on de los
$(k{+}1)$-s\'implices que contienen a $\sigma$,
\[
  \mathbf{m}_\sigma^{\mathrm{up}}
    = \bigoplus_{\rho \in \mathcal{B}^{\mathrm{up}}(\sigma)}
      M^{\mathrm{up}}\bigl(\mathbf{h}_\sigma, \mathbf{h}_\rho\bigr),
\]
donde $\mathcal{B}^{\mathrm{up}}(\sigma) =
\{\rho : [\partial_{k+1} \partial_{k+1}^{\!\top}]_{\sigma,\rho} \neq 0\}$ son los
$k$-s\'implices que comparten un $(k{+}1)$-s\'implice.
\item \textbf{Actualizaci\'on}:
$\mathbf{h}_\sigma' = U\bigl(\mathbf{h}_\sigma,
\mathbf{m}_\sigma^{\mathrm{down}}, \mathbf{m}_\sigma^{\mathrm{up}}\bigr)$.
\end{enumerate}
Aqu\'i $\bigoplus$ denota una agregaci\'on permutaci\'on-invariante (suma, media,
m\'aximo), $M^{\mathrm{down}}, M^{\mathrm{up}}$ son funciones de mensaje
aprendibles y $U$ es la funci\'on de actualizaci\'on~\citep{bodnar2021weisfeiler}.
\end{definition}

\begin{proposition}[El TDL como instancia generalizada del marco MPNN]
\label{prop:priors:tdl-mpnn}
El paso de mensajes simplicial (\Cref{def:priors:simplicial-mp}) es una instancia
generalizada del marco MPNN (\Cref{def:priors:mpnn}) con dominio el complejo
simplicial $\mathcal{K}$, los $k$-s\'implices $\sigma \in \mathcal{K}_k$ como
``nodos''; vecindario doble
$\Nb(\sigma) = \mathcal{B}^{\mathrm{down}}(\sigma) \cup
\mathcal{B}^{\mathrm{up}}(\sigma)$ determinado por los Laplacianos inferior
$L_k^{\mathrm{down}} = \partial_k^{\!\top} \partial_k$ y superior
$L_k^{\mathrm{up}} = \partial_{k+1} \partial_{k+1}^{\!\top}$; la combinaci\'on de
los mensajes inferior y superior; y la actualizaci\'on aprendible
$U(\mathbf{h}_\sigma, \mathbf{m}_\sigma^{\mathrm{down}},
\mathbf{m}_\sigma^{\mathrm{up}})$. En particular, para $k = 0$ se recupera la MPNN
est\'andar sobre grafos.
\end{proposition}

\begin{proof}
Verificamos la correspondencia con las tres componentes del marco MPNN. Fijemos
$k \ge 0$ y consideremos el grafo $\mathcal{G}_k = (\mathcal{K}_k, E_k)$ cuyos
nodos son los $k$-s\'implices y cuyas aristas codifican la adyacencia dada por
$L_k$.

\emph{(1) Mensaje.} Para cada $k$-s\'implice $\sigma$, el paso de mensajes
simplicial calcula
\[
  \mathbf{m}_\sigma^{\mathrm{down}}
    = \bigoplus_{\tau \in \mathcal{B}^{\mathrm{down}}(\sigma)}
      M^{\mathrm{down}}(\mathbf{h}_\sigma, \mathbf{h}_\tau),
  \qquad
  \mathbf{m}_\sigma^{\mathrm{up}}
    = \bigoplus_{\rho \in \mathcal{B}^{\mathrm{up}}(\sigma)}
      M^{\mathrm{up}}(\mathbf{h}_\sigma, \mathbf{h}_\rho).
\]
Definiendo $\Nb(\sigma) = \mathcal{B}^{\mathrm{down}}(\sigma) \cup
\mathcal{B}^{\mathrm{up}}(\sigma)$ y el mensaje conjunto
$\widetilde{M}(\mathbf{h}_\sigma, \mathbf{h}_\nu) =
M^{\mathrm{down}}(\mathbf{h}_\sigma, \mathbf{h}_\nu)$ si
$\nu \in \mathcal{B}^{\mathrm{down}}(\sigma)$ y
$\widetilde{M}(\mathbf{h}_\sigma, \mathbf{h}_\nu) =
M^{\mathrm{up}}(\mathbf{h}_\sigma, \mathbf{h}_\nu)$ si
$\nu \in \mathcal{B}^{\mathrm{up}}(\sigma)$, esto adopta la forma
$m_\sigma = \bigoplus_{\nu \in \Nb(\sigma)} \widetilde{M}(\mathbf{h}_\sigma,
\mathbf{h}_\nu)$, que es exactamente la fase de mensaje de
\Cref{def:priors:mpnn}.

\emph{(2) Agregaci\'on.} El operador $\bigoplus$ (suma, media o m\'aximo)
desempe\~na el papel de la agregaci\'on permutaci\'on-invariante del esquema MPNN.

\emph{(3) Actualizaci\'on.} La funci\'on
$U(\mathbf{h}_\sigma, \mathbf{m}_\sigma^{\mathrm{down}},
\mathbf{m}_\sigma^{\mathrm{up}})$ corresponde directamente a la funci\'on de
actualizaci\'on $U_\ell$ del marco MPNN.

\emph{Reducci\'on a $k = 0$.} Para $0$-s\'implices (nodos), $\partial_0 = 0$
implica $L_0^{\mathrm{down}} = 0$, por lo que $\mathcal{B}^{\mathrm{down}}(v) =
\emptyset$ y el mensaje inferior se anula. El mensaje superior se reduce a
$\mathbf{m}_v^{\mathrm{up}} = \bigoplus_{u \in \mathcal{B}^{\mathrm{up}}(v)}
M^{\mathrm{up}}(\mathbf{h}_v, \mathbf{h}_u)$, donde
$\mathcal{B}^{\mathrm{up}}(v) = \{u : [\partial_1 \partial_1^{\!\top}]_{vu} \neq 0\}
= \Nb(v)$ es el vecindario est\'andar de $v$ en el grafo $1$-esqueleto. Esto
recupera exactamente el paso de mensajes cl\'asico sobre grafos. \qedhere
\end{proof}

\begin{proposition}[Expresividad del paso de mensajes simplicial]
\label{prop:priors:simplicial-expressivity}
El paso de mensajes simplicial es estrictamente m\'as expresivo que el paso de
mensajes en grafos. Existen pares de complejos simpliciales no isomorfos cuyos
$1$-esqueletos son id\'enticos (y por tanto indistinguibles por MPNN) pero que las
redes simpliciales distinguen correctamente gracias a la informaci\'on
topol\'ogica de orden superior.
\end{proposition}

Este resultado, debido a \citet{bodnar2021weisfeiler}, muestra que la
informaci\'on topol\'ogica codificada en los operadores de frontera $\partial_k$
proporciona un poder discriminativo que supera las limitaciones del test de
Weisfeiler--Leman establecidas en la \Cref{sec:priors:wl}. Cierra el arco de
unificaci\'on del marco MPNN: mientras que las ANN, las CNN y los
\emph{transformers} codifican sesgos inductivos espaciales o relacionales sobre
nodos, el TDL introduce un \emph{sesgo inductivo topol\'ogico de orden superior}
que captura interacciones entre tri\'angulos, tetraedros y $k$-s\'implices
arbitrarios, proporcionando un poder expresivo estrictamente superior y superando
la barrera de Weisfeiler--Leman~\citep{bodnar2021weisfeiler,
hajij2023topological}.

\section{Poder expresivo y sus l\'imites}
\label{sec:priors:expressivity-limits}

El paso de mensajes es unificador pero no omnipotente. Su poder expresivo sobre
grafos est\'a fijado por un algoritmo combinatorio cl\'asico, y su profundidad
est\'a limitada por dos patolog\'ias complementarias. Las desarrollamos aqu\'i; el
an\'alisis espectral y geom\'etrico reaparece en \Cref{ch:graphs,ch:geodesics}.

\subsection{El test de Weisfeiler--Leman}
\label{sec:priors:wl}

El \emph{test de Weisfeiler--Leman} (1-WL) es un algoritmo cl\'asico, propuesto
originalmente en 1968~\citep{weisfeiler1968reduction} para decidir si dos grafos
son isomorfos mediante un proceso iterativo de \emph{refinamiento de colores}. Su
importancia en GDL radica en que establece una cota superior fundamental para el
poder expresivo de las MPNN.

\begin{definition}[Algoritmo de Weisfeiler--Leman (1-WL)]
\label{def:priors:wl}
Sea $\mathcal{G} = (V, E)$ un grafo con $n$ v\'ertices. El \emph{algoritmo de
Weisfeiler--Leman de orden 1} (1-WL) consiste en el siguiente proceso iterativo
de coloraci\'on.
\begin{enumerate}
\item \textbf{Inicializaci\'on} ($t = 0$): asignar a cada nodo $v \in V$ un color
inicial $c^{(0)}(v)$, t\'ipicamente id\'entico para todos los nodos (o basado en
etiquetas de nodos si las hay).
\item \textbf{Refinamiento iterativo}: en cada iteraci\'on $t \ge 0$, actualizar
el color de cada nodo seg\'un
\[
  c^{(t+1)}(v) = \HASH\Bigl(c^{(t)}(v),\,
      \{\!\{ c^{(t)}(u) : u \in \Nb(v) \}\!\}\Bigr),
\]
donde $\Nb(v)$ es el vecindario de $v$, $\{\!\{\cdot\}\!\}$ denota un
multiconjunto y $\HASH$ es una funci\'on inyectiva que asigna un nuevo color
\'unico a cada par distinto (color antiguo, multiconjunto de colores vecinos).
\item \textbf{Convergencia}: el algoritmo termina cuando la partici\'on de colores
se estabiliza, es decir, cuando $c^{(t+1)}(v) = c^{(t)}(v)$ para todo $v \in V$.
Esto ocurre en a lo sumo $n$ iteraciones.
\item \textbf{Histograma final}: el resultado es el \emph{histograma de colores}
$\mathbf{h}(\mathcal{G}) = (|\{v : c^{(T)}(v) = k\}|)_k$, que cuenta cu\'antos
nodos portan cada color final.
\end{enumerate}
\end{definition}

\begin{definition}[WL-equivalencia]
\label{def:priors:wl-equivalence}
Dos grafos $\mathcal{G}_1$ y $\mathcal{G}_2$ son \emph{WL-equivalentes}, denotado
$\mathcal{G}_1 \equiv_{\mathrm{WL}} \mathcal{G}_2$, si el algoritmo 1-WL produce
el mismo histograma de colores para ambos.
\end{definition}

\begin{figure}[h]
\centering
\begin{tikzpicture}[>=Stealth, font=\small,
    wlnode/.style={circle, draw, thick, minimum size=7mm, inner sep=0pt},
    stage/.style={font=\bfseries\small}
]
\newcommand{\drawedges}{%
    \draw[thick, gdlgray!50] (n1) -- (n2);
    \draw[thick, gdlgray!50] (n2) -- (n3);
    \draw[thick, gdlgray!50] (n3) -- (n4);
    \draw[thick, gdlgray!50] (n4) -- (n1);
    \draw[thick, gdlgray!50] (n3) -- (n5);
    \draw[thick, gdlgray!50] (n4) -- (n5);
    \draw[thick, gdlgray!50] (n1) -- (n6);
    \draw[thick, gdlgray!50] (n2) -- (n6);
    \draw[thick, gdlgray!50] (n3) -- (n6);
    \draw[thick, gdlgray!50] (n4) -- (n6);
}
% (a) inicial
\begin{scope}[shift={(0,0)}]
    \node[stage] at (1, 3) {(a) Inicial: $c^{(0)}$};
    \node[wlnode, fill=gdlgray] (n1) at (0, 1.8) {};
    \node[wlnode, fill=gdlgray] (n2) at (2, 1.8) {};
    \node[wlnode, fill=gdlgray] (n3) at (2, 0) {};
    \node[wlnode, fill=gdlgray] (n4) at (0, 0) {};
    \node[wlnode, fill=gdlgray] (n5) at (1, -1.2) {};
    \node[wlnode, fill=gdlgray] (n6) at (1, 0.9) {};
    \drawedges
    \node[font=\tiny, above left=-1pt] at (n1) {$v_1$};
    \node[font=\tiny, above right=-1pt] at (n2) {$v_2$};
    \node[font=\tiny, right=1pt] at (n3) {$v_3$};
    \node[font=\tiny, left=1pt] at (n4) {$v_4$};
    \node[font=\tiny, below=1pt] at (n5) {$v_5$};
    \node[font=\tiny, above=1pt] at (n6) {$v_6$};
\end{scope}
\draw[->, very thick, gdlgray!50] (3.3, 0.9) -- node[above, font=\scriptsize] {HASH} (4.7, 0.9);
% (b) refinamiento por grado
\begin{scope}[shift={(5.5,0)}]
    \node[stage] at (1, 3) {(b) Iteraci\'on 1: $c^{(1)}$};
    \node[wlnode, fill=gdlblue!50] (n1) at (0, 1.8) {};
    \node[wlnode, fill=gdlblue!50] (n2) at (2, 1.8) {};
    \node[wlnode, fill=gdlred!50] (n3) at (2, 0) {};
    \node[wlnode, fill=gdlred!50] (n4) at (0, 0) {};
    \node[wlnode, fill=gdlgreen!50] (n5) at (1, -1.2) {};
    \node[wlnode, fill=gdlred!50] (n6) at (1, 0.9) {};
    \drawedges
\end{scope}
\draw[->, very thick, gdlgray!50] (8.8, 0.9) -- node[above, font=\scriptsize] {HASH} (10.2, 0.9);
% (c) refinamiento adicional
\begin{scope}[shift={(11,0)}]
    \node[stage] at (1, 3) {(c) Iteraci\'on 2: $c^{(2)}$};
    \node[wlnode, fill=gdlblue!50] (n1) at (0, 1.8) {};
    \node[wlnode, fill=gdlblue!50] (n2) at (2, 1.8) {};
    \node[wlnode, fill=gdlyellow!60] (n3) at (2, 0) {};
    \node[wlnode, fill=gdlyellow!60] (n4) at (0, 0) {};
    \node[wlnode, fill=gdlgreen!50] (n5) at (1, -1.2) {};
    \node[wlnode, fill=gdlred!50] (n6) at (1, 0.9) {};
    \drawedges
    \node[font=\scriptsize, text=black!50, anchor=west] at (2.3, 0.9) {estable};
\end{scope}
\node[anchor=west, font=\scriptsize] at (0, -2.6) {%
    Refinamiento: $c^{(t+1)}(v) = \mathrm{HASH}\!\big(c^{(t)}(v),\, \{\!\{c^{(t)}(u) : u \in \mathcal{N}(v)\}\!\}\big)$};
\node[anchor=west, font=\scriptsize, text=black!60] at (0, -3.3) {%
    Cada nodo actualiza su color seg\'un su color actual y el multiconjunto de colores de sus vecinos.};
\node[anchor=west] at (0, -4.2) {%
    \tikz{\node[wlnode, fill=gdlblue!50, minimum size=4mm] {};} {\scriptsize grado 3} \quad
    \tikz{\node[wlnode, fill=gdlred!50, minimum size=4mm] {};} {\scriptsize grado 4 (tipo A)} \quad
    \tikz{\node[wlnode, fill=gdlyellow!60, minimum size=4mm] {};} {\scriptsize grado 4 (tipo B)} \quad
    \tikz{\node[wlnode, fill=gdlgreen!50, minimum size=4mm] {};} {\scriptsize grado 2}
};
\end{tikzpicture}
\caption[Refinamiento de colores Weisfeiler--Leman]{El test de refinamiento de
colores 1-WL sobre un grafo de seis nodos. (a) Coloraci\'on inicial uniforme.
(b) Tras la primera iteraci\'on, los nodos se distinguen por grado: azul (grado
3), rojo (grado 4), verde (grado 2). (c) La segunda iteraci\'on refina seg\'un el
multiconjunto de colores vecinos: $v_3$ y $v_4$ (amarillo) se separan de $v_6$
(rojo) porque sus vecindarios tienen composiciones de colores distintas.}
\label{fig:priors:wl}
\end{figure}

La conexi\'on profunda entre el test WL y las redes de grafos fue establecida
rigurosamente por \citet{xu2019how} y \citet{morris2019weisfeiler}.

\begin{theorem}[Equivalencia WL--MPNN]
\label{thm:priors:wl-mpnn}
Sea $\mathcal{G}$ una clase de grafos y sea $\mathcal{A}$ una MPNN con funciones
de agregaci\'on y actualizaci\'on inyectivas. Entonces:
\begin{enumerate}
\item \textbf{Cota superior}: $\mathcal{A}$ no puede distinguir grafos
WL-equivalentes. Es decir, si
$\mathcal{G}_1 \equiv_{\mathrm{WL}} \mathcal{G}_2$, entonces
$\mathcal{A}(\mathcal{G}_1) = \mathcal{A}(\mathcal{G}_2)$.
\item \textbf{Optimalidad}: existe una MPNN (espec\'ificamente, una Graph
Isomorphism Network) que es \emph{tan expresiva como} el test WL. Es decir, si
$\mathcal{G}_1 \not\equiv_{\mathrm{WL}} \mathcal{G}_2$, existe una configuraci\'on
de pesos tal que
$\mathcal{A}(\mathcal{G}_1) \neq \mathcal{A}(\mathcal{G}_2)$.
\end{enumerate}
\end{theorem}

\begin{proof}[Esbozo de demostraci\'on]
La clave es la correspondencia estructural entre una iteraci\'on de WL y una capa
de una MPNN.

\emph{WL $\to$ MPNN}: el refinamiento de WL
$c^{(t+1)}(v) = \HASH(c^{(t)}(v), \{\!\{c^{(t)}(u)\}\!\})$ se mapea directamente a
una capa mensaje-agregaci\'on-actualizaci\'on
\[
  h_v^{(t+1)} = \phi\Bigl(h_v^{(t)},
      \bigoplus_{u \in \Nb(v)} \psi(h_u^{(t)})\Bigr).
\]
Si $\phi$ y $\bigoplus$ son suficientemente expresivas (por ejemplo, una suma
sobre \emph{embeddings} \'unicos seguida de un MLP), reproducen exactamente el
comportamiento de $\HASH$.

\emph{MPNN $\to$ WL}: rec\'iprocamente, si dos grafos son WL-equivalentes
comparten la misma secuencia de particiones de colores. Dado que una MPNN solo
puede agregar informaci\'on del vecindario, no puede ``romper'' la simetr\'ia que
el test WL preserva.

La demostraci\'on completa requiere mostrar que el agregador de suma sobre
\emph{embeddings} suficientemente \'unicos es inyectivo, lo cual se sigue del
teorema de aproximaci\'on universal para conjuntos~\citep{zaheer2017deepsets}.
\qedhere
\end{proof}

\begin{example}[Grafos no isomorfos indistinguibles por WL]
\label{ex:priors:wl-indistinguishable}
Los grafos no isomorfos m\'as simples que el test 1-WL no puede distinguir son los
grafos regulares de igual grado y tama\~no. Consideremos:
\begin{itemize}
\item $\mathcal{G}_1$: un ciclo de seis nodos (hex\'agono), $C_6$;
\item $\mathcal{G}_2$: dos tri\'angulos disjuntos, $K_3 \sqcup K_3$.
\end{itemize}
Ambos grafos tienen $6$ nodos, $6$ aristas y todos los nodos tienen grado $2$. El
algoritmo WL asigna:
\begin{enumerate}
\item Iteraci\'on 0: todos los nodos reciben el mismo color inicial
$c^{(0)} = 1$.
\item Iteraci\'on 1: cada nodo ve vecinos de color $1$, multiconjunto
$\{\!\{1, 1\}\!\}$. Nuevo color $c^{(1)} = \HASH(1, \{\!\{1, 1\}\!\}) = 2$ para
todos.
\item Convergencia: la partici\'on no cambia. Histograma final
$\mathbf{h}(\mathcal{G}_1) = \mathbf{h}(\mathcal{G}_2) = (6)$.
\end{enumerate}
Sin embargo, $\mathcal{G}_1 \not\cong \mathcal{G}_2$: el ciclo $C_6$ es conexo
mientras que $K_3 \sqcup K_3$ tiene dos componentes conexas. Esta limitaci\'on
muestra que las MPNN est\'andar no pueden detectar la conectividad global de
grafos regulares. Si una tarea requiere distinguir mol\'eculas con estructura
global similar pero conectividad diferente (por ejemplo, is\'omeros
estructurales), las MPNN est\'andar pueden fallar sistem\'aticamente.
\end{example}

\subsection{Sobre-suavizado: an\'alisis espectral}
\label{sec:priors:oversmoothing}

La \Cref{thm:priors:wl-mpnn} delimita qu\'e grafos pueden distinguir las MPNN,
pero no dice nada sobre qu\'e sucede cuando apilamos muchas capas de paso de
mensajes. La respuesta, conocida como \emph{sobre-suavizado}
(\emph{oversmoothing})~\citep{li2018deeper}, es que las representaciones de los
nodos convergen a un valor uniforme, perdiendo todo poder discriminativo. Este
fen\'omeno admite una explicaci\'on espectral elegante.

\begin{definition}[Energ\'ia de Dirichlet en grafos]
\label{def:priors:dirichlet-energy}
Sea $\mathcal{G} = (V, E)$ un grafo con Laplaciano $L_G$ y sea
$H \in \R^{n \times d}$ una matriz de caracter\'isticas de nodos ($n = |V|$, $d$
canales). La \emph{energ\'ia de Dirichlet} de $H$ es
\[
  E(H) = \tr(H^{\!\top} L_G H)
       = \sum_{c=1}^{d} \mathbf{h}_c^{\!\top} L_G \mathbf{h}_c
       = \sum_{c=1}^{d} \sum_{(i,j) \in E} W_{ij}\,(h_{ic} - h_{jc})^2,
\]
donde $\mathbf{h}_c$ denota la $c$-\'esima columna de $H$.
\end{definition}

La energ\'ia de Dirichlet mide la variaci\'on total de las se\~nales $H$ sobre el
grafo: $E(H) = 0$ si y solo si $H$ es constante en cada componente conexa.

\begin{proposition}[MPNN lineal como difusi\'on discreta]
\label{prop:priors:mpnn-diffusion}
Consideremos una MPNN (\Cref{def:priors:mpnn}) con agregaci\'on por media y sin no
linealidad,
\[
  h_v^{(\ell+1)} = (1 - \alpha)\, h_v^{(\ell)}
      + \alpha \sum_{u \in \Nb(v)} \frac{h_u^{(\ell)}}{\deg(v)},
  \qquad \alpha \in (0, 1].
\]
Sea $\tilde L = D^{-1} L_G = I - D^{-1} A$ el Laplaciano por paseo aleatorio, que
comparte su espectro con el Laplaciano normalizado sim\'etrico. Entonces la
actualizaci\'on se escribe matricialmente como
\[
  H^{(\ell+1)} = (I - \alpha \tilde L)\, H^{(\ell)}.
\]
\end{proposition}

\begin{proof}
Para cada nodo $v$ la regla de actualizaci\'on es
\[
  h_v^{(\ell+1)}
    = (1 - \alpha) h_v^{(\ell)} + \alpha \sum_{u \sim v} \frac{h_u^{(\ell)}}{\deg(v)}
    = h_v^{(\ell)} - \alpha\Bigl(h_v^{(\ell)} - \sum_{u \sim v} \frac{h_u^{(\ell)}}{\deg(v)}\Bigr)
    = h_v^{(\ell)} - \alpha (\tilde L\, H^{(\ell)})_v,
\]
usando $\tilde L = D^{-1} L_G = I - D^{-1} A$ y
$(D^{-1} A H)_v = \sum_{u \sim v} h_u / \deg(v)$. \qedhere
\end{proof}

Esta es la discretizaci\'on temporal de la ecuaci\'on del calor
$\partial_t H = -\alpha \tilde L H$ en el grafo. Iterando $\ell$ capas se obtiene
$H^{(\ell)} = (I - \alpha \tilde L)^\ell H^{(0)}$, que difunde progresivamente las
se\~nales.

\begin{theorem}[Convergencia exponencial del sobre-suavizado]
\label{thm:priors:oversmoothing}
Sea $\mathcal{G}$ un grafo conexo con valores propios del Laplaciano normalizado
$0 = \lambda_1 < \lambda_2 \le \cdots \le \lambda_n$ y
$\alpha \in (0, 1/\lambda_n)$. Entonces la energ\'ia de Dirichlet de las
representaciones tras $\ell$ capas de la MPNN lineal satisface
\[
  E\bigl(H^{(\ell)}\bigr) \le (1 - \alpha \lambda_2)^{2\ell}\, E\bigl(H^{(0)}\bigr).
\]
En particular, $E(H^{(\ell)}) \to 0$ exponencialmente cuando $\ell \to \infty$:
todas las representaciones convergen a un mismo vector.
\end{theorem}

\begin{proof}
Sea $\{u_k\}_{k=1}^{n}$ la base ortonormal de autovectores de
$L_{\mathrm{sym}} = D^{-1/2} L_G D^{-1/2}$ con autovalores $\lambda_k$. Para cada
columna definimos $\mathbf{f}_c = D^{1/2} \mathbf{h}_c$, de modo que
$\mathbf{h}_c^{\!\top} L_G \mathbf{h}_c = \mathbf{f}_c^{\!\top} L_{\mathrm{sym}}
\mathbf{f}_c$. La transformada de Fourier en grafos descompone
$\mathbf{f}_c = \sum_{k=1}^{n} \hat f_{c,k}\, u_k$.

El operador de difusi\'on satisface
$D^{1/2}(I - \alpha \tilde L) D^{-1/2} = I - \alpha L_{\mathrm{sym}}$, con
autovalores $1 - \alpha \lambda_k$. Tras $\ell$ iteraciones, los coeficientes
espectrales de $\mathbf{f}_c^{(\ell)} = D^{1/2} \mathbf{h}_c^{(\ell)}$ se escalan
como $(1 - \alpha \lambda_k)^\ell \hat f_{c,k}$. Por tanto, la energ\'ia de
Dirichlet en el dominio espectral es
\[
  E(H^{(\ell)})
    = \sum_{c=1}^{d} \mathbf{f}_c^{(\ell)\top} L_{\mathrm{sym}} \mathbf{f}_c^{(\ell)}
    = \sum_{c=1}^{d} \sum_{k=2}^{n} \lambda_k (1 - \alpha \lambda_k)^{2\ell}
      |\hat f_{c,k}|^2,
\]
donde el t\'ermino $k = 1$ desaparece porque $\lambda_1 = 0$. Como
$\alpha < 1/\lambda_n$, se tiene $|1 - \alpha \lambda_k| < 1$ para $k \ge 2$, y en
particular $|1 - \alpha \lambda_k| \le 1 - \alpha \lambda_2$ para todo $k \ge 2$
(el m\'aximo se alcanza en $\lambda_2$, la frecuencia m\'as baja no trivial). Por
tanto,
\[
  E(H^{(\ell)})
    \le (1 - \alpha \lambda_2)^{2\ell} \sum_{c=1}^{d} \sum_{k=2}^{n}
        \lambda_k |\hat f_{c,k}|^2
    = (1 - \alpha \lambda_2)^{2\ell}\, E(H^{(0)}).
\]
Como $0 < \alpha \lambda_2 < 1$, el factor $(1 - \alpha \lambda_2)^{2\ell} \to 0$
exponencialmente. \qedhere
\end{proof}

\begin{remark}[Interpretaci\'on geom\'etrica del sobre-suavizado]
\label{rem:priors:oversmoothing}
La tasa de convergencia $(1 - \alpha \lambda_2)^{2\ell}$ est\'a controlada por el
\emph{gap espectral} $\lambda_2$, el segundo autovalor del Laplaciano. Los grafos
con $\lambda_2$ grande (como los expansores) convergen m\'as r\'apido al estado
uniforme, mientras que los grafos poco conectados ($\lambda_2 \approx 0$)
preservan mejor la diversidad de representaciones. Esta es la contrapartida
discreta exacta de la difusi\'on del calor en variedades riemannianas, donde el
gap espectral del operador de Laplace--Beltrami controla la tasa de mezclado.
Entre las estrategias para mitigar el sobre-suavizado destacan las conexiones
residuales (que preservan la se\~nal original), la normalizaci\'on por capas y los
enfoques topol\'ogicos como Neural Sheaf Diffusion~\citep{bodnar2022neural}, donde
se reemplaza la difusi\'on escalar por una difusi\'on en haces (\emph{sheaves})
que permite a cada arista transportar informaci\'on en espacios vectoriales
distintos.
\end{remark}

\subsection{Sobre-compresi\'on: cuellos de botella topol\'ogicos}
\label{sec:priors:oversquashing}

Si el sobre-suavizado surge de la p\'erdida de se\~nal de alta frecuencia, el
fen\'omeno complementario de \emph{sobre-compresi\'on} (\emph{oversquashing}) es
la incapacidad de propagar informaci\'on a larga distancia. Intuitivamente, la
informaci\'on de un nodo lejano debe comprimirse exponencialmente para ``pasar''
por los cuellos de botella del grafo.

\begin{definition}[Sensibilidad jacobiana de una MPNN]
\label{def:priors:sensitivity}
Para una MPNN con representaciones $h_v^{(\ell)} \in \R^d$, la
\emph{sensibilidad} del nodo $v$ respecto al nodo $u$ tras $\ell$ capas es
\[
  S_{vu}^{(\ell)} = \left\| \frac{\partial h_v^{(\ell)}}{\partial h_u^{(0)}} \right\|,
\]
la norma de operador de la matriz jacobiana
$\partial h_v^{(\ell)} / \partial h_u^{(0)} \in \R^{d \times d}$.
\end{definition}

Si $S_{vu}^{(\ell)} \approx 0$, el nodo $v$ es esencialmente insensible a la
entrada del nodo $u$: cualquier informaci\'on que $u$ porte se ha perdido durante
la propagaci\'on.

\begin{proposition}[Decaimiento exponencial de la sensibilidad]
\label{prop:priors:oversquashing}
Sea una MPNN con funciones de mensaje y actualizaci\'on Lipschitz con constantes
$c_M$ y $c_U$ respectivamente, y sea $\hat A = D^{-1} A$ la matriz de transici\'on
del paseo aleatorio en $\mathcal{G}$. Entonces, para cualesquiera nodos
$v, u \in V$,
\[
  S_{vu}^{(\ell)} \le (c_U \cdot c_M)^\ell\, [\hat A^\ell]_{vu}.
\]
En particular, si la distancia en el grafo satisface
$d_{\mathcal{G}}(v, u) > \ell$, entonces $[\hat A^\ell]_{vu} = 0$ y la
sensibilidad es exactamente cero.
\end{proposition}

\begin{proof}[Esbozo de demostraci\'on]
Procedemos por inducci\'on sobre $\ell$. Para $\ell = 1$ y nodos $v \neq u$, como
$h_v^{(1)} = U(h_v^{(0)}, m_v^{(1)})$, la regla de la cadena da
\[
  \frac{\partial h_v^{(1)}}{\partial h_u^{(0)}}
    = \frac{\partial U}{\partial h_v^{(0)}}
      \cdot \underbrace{\frac{\partial h_v^{(0)}}{\partial h_u^{(0)}}}_{=\,0 \text{ si } v \neq u}
    + \frac{\partial U}{\partial m_v} \cdot \frac{\partial m_v}{\partial h_u^{(0)}}
    = \frac{\partial U}{\partial m_v} \cdot \frac{\partial m_v}{\partial h_u^{(0)}}.
\]
El segundo factor es no nulo solo si $u \in \Nb(v)$, en cuyo caso la constante de
Lipschitz de $M$ aporta $c_M$ y la normalizaci\'on por grado aporta
$1/\deg(v) = [\hat A]_{vu}$. As\'i, $S_{vu}^{(1)} \le c_U c_M [\hat A]_{vu}$.

Para el paso inductivo, la regla de la cadena multivariable da
\[
  \frac{\partial h_v^{(\ell+1)}}{\partial h_u^{(0)}}
    = \sum_{w \in \Nb(v) \cup \{v\}}
      \frac{\partial h_v^{(\ell+1)}}{\partial h_w^{(\ell)}}
      \cdot \frac{\partial h_w^{(\ell)}}{\partial h_u^{(0)}}.
\]
Acotando cada factor por las constantes de Lipschitz y usando la hip\'otesis
inductiva para $\partial h_w^{(\ell)} / \partial h_u^{(0)}$, obtenemos el producto
matricial $\hat A^{\ell+1}$ que registra todos los caminos de longitud $\ell + 1$
de $u$ a $v$. \qedhere
\end{proof}

\begin{remark}[Curvatura y sobre-compresi\'on]
\label{rem:priors:curvature}
La cota revela que la topolog\'ia del grafo, codificada en $\hat A^\ell$, es el
factor determinante de la sobre-compresi\'on. \citet{topping2022understanding}
formalizan esta conexi\'on mediante la \emph{curvatura de Ricci discreta}
(Ollivier): aristas $(u, v)$ con curvatura de Ricci negativa corresponden a
cuellos de botella donde $[\hat A^\ell]_{vu}$ decae r\'apidamente. Esta es la
contrapartida discreta de la ecuaci\'on de Jacobi en geometr\'ia riemanniana: la
curvatura seccional negativa hace que las geod\'esicas diverjan, reduciendo el
``volumen de informaci\'on'' que puede propagarse. Topping et al.\ proponen el
\emph{reconfigurado} (\emph{rewiring}) del grafo mediante flujo de Ricci,
a\~nadiendo aristas en regiones de curvatura negativa para aliviar los cuellos de
botella ---un ejemplo de c\'omo la geometr\'ia diferencial prescribe soluciones a
problemas de las GNN.
\end{remark}

\subsection{La jerarqu\'ia k-WL y extensiones}
\label{sec:priors:kwl}

La \Cref{thm:priors:wl-mpnn} establece que las MPNN son tan expresivas como el
test 1-WL, que opera sobre nodos individuales. Una extensi\'on natural considera
tests de orden superior que operan sobre tuplas de nodos.

\begin{definition}[Test de Weisfeiler--Leman de orden $k$]
\label{def:priors:kwl}
El \emph{test $k$-WL} opera sobre $k$-tuplas de nodos
$\bar v = (v_1, \dots, v_k) \in V^k$~\citep{morris2019weisfeiler}. Se define una
coloraci\'on inicial $c^{(0)}(\bar v)$ basada en el tipo de isomorfismo de la
subtupla, y el refinamiento iterativo
\[
  c^{(t+1)}(\bar v) = \HASH\Bigl(c^{(t)}(\bar v),\,
    \bigl\{\!\!\bigl\{ \bigl(c^{(t)}(\bar v[w/1]), \dots,
      c^{(t)}(\bar v[w/k])\bigr) : w \in V \bigr\}\!\!\bigr\}\Bigr),
\]
donde $\bar v[w/i]$ denota la tupla obtenida al reemplazar la $i$-\'esima
componente de $\bar v$ por $w$.
\end{definition}

\begin{proposition}[Jerarqu\'ia estricta de expresividad]
\label{prop:priors:kwl-hierarchy}
Los tests de Weisfeiler--Leman forman una jerarqu\'ia estrictamente creciente en
poder expresivo,
\[
  1\text{-WL} \subsetneq 2\text{-WL} \subsetneq 3\text{-WL} \subsetneq \cdots,
\]
es decir, para cada $k \ge 1$ existen grafos que $k$-WL no distingue pero
$(k+1)$-WL s\'i. La complejidad computacional del test $k$-WL es
$O(n^k \log n)$ por iteraci\'on, donde $n = |V|$.
\end{proposition}

\begin{proof}
La demostraci\'on de la separaci\'on estricta requiere t\'ecnicas de la teor\'ia
de modelos finitos (en particular, la conexi\'on con la l\'ogica $C^k$ de
Immerman--Lander) y excede el alcance de este libro; v\'ease
\citet{morris2019weisfeiler} para los detalles. La complejidad $O(n^k \log n)$ se
sigue directamente del refinamiento iterativo: hay $n^k$ tuplas, cada iteraci\'on
inspecciona $n$ sustituciones por tupla, y la estabilizaci\'on ocurre en a lo sumo
$O(n^k)$ iteraciones con coste $O(\log n)$ del hash. \qedhere
\end{proof}

Existen diversas estrategias para superar la barrera 1-WL sin incurrir en el coste
exponencial de $k$-WL:
\begin{itemize}
\item \textbf{Paso de mensajes simplicial}: como muestra la
\Cref{prop:priors:simplicial-expressivity}, operar sobre complejos simpliciales
($k$-cadenas) supera 1-WL incorporando informaci\'on topol\'ogica de orden
superior.
\item \textbf{Codificaciones estructurales}: a\~nadir codificaciones posicionales
basadas en autovectores del Laplaciano o en paseos aleatorios rompe la simetr\'ia
de los nodos.
\item \textbf{Graph transformers espectrales}: arquitecturas como la de
\citet{kreuzer2021spectral} combinan atenci\'on global con codificaciones
espectrales, conectando el an\'alisis espectral con los mecanismos de atenci\'on.
\end{itemize}

Las tres limitaciones analizadas (sobre-suavizado, sobre-compresi\'on y barrera
WL) son manifestaciones del car\'acter fundamentalmente \emph{local} del paso de
mensajes: cada capa propaga informaci\'on un solo salto. El sobre-suavizado
muestra que demasiadas capas destruyen la se\~nal discriminativa; la
sobre-compresi\'on muestra que pocas capas impiden la comunicaci\'on a larga
distancia; y la barrera WL muestra que ciertos patrones globales son indetectables
localmente. Las herramientas geom\'etricas desarrolladas en este libro (an\'alisis
espectral, curvatura discreta, topolog\'ia de orden superior) no solo
\emph{diagnostican} estas limitaciones: \emph{prescriben} soluciones concretas.

\section{Resumen}
\label{sec:priors:summary}

La invarianza y la equivarianza convierten la exigencia informal de que un modelo
``respete las simetr\'ias de los datos'' en restricciones precisas sobre una
funci\'on. Como las aplicaciones equivariantes se componen
(\Cref{thm:priors:composition}), estas restricciones producen un esquema
constructivo: apilar capas equivariantes y engrosamiento local, y terminar con
una lectura invariante. El teorema de aproximaci\'on universal equivariante
(\Cref{thm:priors:universal}) certifica que esta disciplina no cuesta poder
expresivo. El paso de mensajes hace concreto el esquema y revela las
arquitecturas completamente conectadas, convolucionales, de atenci\'on, de grafos
e incluso simpliciales como una sola construcci\'on bajo distintos sesgos
geom\'etricos, con la jerarqu\'ia WL, el sobre-suavizado y la sobre-compresi\'on
delimitando su alcance expresivo. La perspectiva unificadora del GDL expone una
progresi\'on natural en sofisticaci\'on geom\'etrica: las redes completamente
conectadas operan en espacios euclidianos con simetr\'ias triviales; las
convoluciones introducen un sesgo inductivo espacial mediante localidad y
equivarianza traslacional; los \emph{transformers} generalizan esto mediante
conectividad global adaptativa; y el paradigma se extiende naturalmente a grafos
arbitrarios y a topolog\'ias de orden superior. Este marco es predictivo, no
meramente descriptivo: sugiere principios de dise\~no para nuevas arquitecturas
h\'ibridas. Los cap\'itulos restantes instancian este esquema sobre cada uno de
los cinco dominios geom\'etricos.
